\documentclass[10pt]{article}
\usepackage[letterpaper,margin=0.68in]{geometry}
\usepackage{amsmath,amssymb,mathtools}
\usepackage{microtype}
\usepackage[hidelinks]{hyperref}
\hypersetup{
  pdftitle={Technical summary: the quantum chromatic number of the G19 join family},
  pdfauthor={Alec Kriebel},
  pdfsubject={Quantum graph theory and quantum chromatic number},
  pdfkeywords={quantum chromatic number, quantum coloring,
    orthogonal representation, quantum graph theory}
}
\setlength{\parindent}{0pt}
\setlength{\parskip}{3pt}
\setlength{\abovedisplayskip}{4pt}
\setlength{\belowdisplayskip}{4pt}
\pagestyle{plain}

\begin{document}
\begin{center}
{\Large\bfseries Technical summary: the $G_{19}$ join family}\\[-1pt]
{\small Alec Kriebel \textbullet\ Independent Researcher \textbullet\ August 1, 2026}
\end{center}
\vspace{-6pt}
\fontsize{10.1}{11.9}\selectfont

\textbf{Result.}
Lalonde's Section~4.2 family is $G_n=G_{19}\vee K_{n-3}$; write it as $J_n$
to avoid overloading the fixed 19-vertex base graph.  The base graph has
graph6 string \texttt{RxLAKA@AgYAWDGO?O?@??A?W@@OC@\_}.  For every
$n\ge3$ and $d,s\ge1$, with $\xi$ denoting complex orthogonal rank,
\[
 \boxed{\xi(J_n)=n<\chi_q(J_n)=n+1,}\qquad
 \chi_q^{[d]}(J_n)=\chi_q^{(s)}(J_n)=n+1.
\]
This confirms Lalonde's candidate finite witnesses for the complex-sphere
obstruction and removes the rank-one restriction from his theorem.  In
particular, for the one-apex graph $H=J_4$, $\chi_q(H)=5$.  Zero projectors,
nonuniform original ranks, arbitrary finite dimension, and noncommuting
projectors are all allowed.

\textbf{Uniformization and clique saturation.}
Suppose $J_n$ has a quantum $n$-coloring in dimension $d$.  Direct-sum over
all color permutations:
\[
 \widetilde P_{v,c}=\bigoplus_{\pi\in S_n}P_{v,\pi(c)}.
\]
The new dimension is $D=n!d=nr$, and every projector has common positive
rank $r=(n-1)!d$.  Thus uniformity is a theorem-level reduction, not an
ansatz.  For an $n$-clique, each fixed-color column consists of pairwise
orthogonal projections.  Its trace is at most $D$; summing over all colors
and using the row partitions makes all $n$ inequalities equalities.  Hence
every clique column sums to $I_D$.

Fix a color $c$.  Let $R_c$ be the sum of its projections on the joined
clique and $Q_c=I-R_c$.  Each of $123$, $189$, $345$, and $267$, joined with
$K_{n-3}$, is an $n$-clique.  Therefore
\[
 \dim Q_c=3r,\qquad \sum_{v\in T_j}P_{v,c}=Q_c.
\]
On $Q_c$, write $e_v=P_{v,c}|_{Q_c}$.  These are rank-$r$ projections giving
a projective orthogonal representation of $G_{19}$.

\textbf{Exact core SOS.}
Set $X=e_{10}+e_{11}+e_{12}+e_{13}$,
\[
\begin{aligned}
H_1&=e_{10}+e_{11}-e_{12}-e_{13},\\
H_2&=e_{10}-e_{11}+e_{12}-e_{13},\\
H_3&=e_{10}-e_{11}-e_{12}+e_{13},
\end{aligned}
\qquad
(d_1,d_2,d_3)=(e_5-e_4,e_7-e_6,e_9-e_8).
\]
Projection, edge, and triangle relations give
\[
 \sum_i d_i^2=2I,\qquad X^2+\sum_iH_i^2=4X,
 \qquad\sum_i\{d_i,H_i\}=-4X.
\]
The final identity is graph-specific: for each $j=10,\ldots,13$, the
positive Walsh signs select exactly the three neighbors of $j$ among
$4,\ldots,9$.  Expanding the following factors now gives the exact rational
identity
\[
 \frac43I-X=\sum_{j=0}^3f_j^*f_j,
 \quad f_0=\frac12\left(X-\frac43I\right),
 \quad f_i=\frac23d_i+\frac12H_i.
\]
Both sides have trace zero, since $\dim Q_c=3r$ and each $e_j$ has rank $r$.
Faithfulness of finite matrix trace forces every $f_j=0$.

\newpage
Put $A=e_4-e_5$, $B=e_6-e_7$, and $C=e_8-e_9$.  Walsh inversion and
$e_j^2=e_j$ yield
\[
 \{B,C\}=A,\qquad\{A,C\}=B,\qquad\{A,B\}=C.
\]
Relative to $Q_c=\operatorname{ran}e_1\oplus\operatorname{ran}e_2
\oplus\operatorname{ran}e_3$, support restrictions put $A,B,C$ on coordinate
pairs $12,13,23$.  The anticommutator equations kill their diagonal blocks;
the remaining blocks are unitaries with one cocycle relation.  A
block-diagonal gauge makes all three blocks $I_r$.  Thus vertices
$1,\ldots,13$ are exactly the scalar sign-vector core tensored with a
multiplicity space $K_c\cong\mathbb C^r$.

\textbf{Exhaustive tail.}
The remaining ranges are encoded by an arbitrary $r$-plane
$M_c\subset K_c\oplus K_c$ invariant under
$J_c=\left(\begin{smallmatrix}0&I\\-I&0\end{smallmatrix}\right)$:
\[
\begin{array}{c|ccc}
v&14&15&16\\ \hline
\operatorname{ran}e_v&(a,0,b)&(0,a,b)&(a,-b,0)
\end{array}\quad (a,b)\in M_c,
\]
\[
\begin{array}{c|ccc}
v&17&18&19\\ \hline
\operatorname{ran}e_v&(x,x,z)&(x,-z,z)&(x,-z,x)
\end{array}\quad (x,z)\in M_c^\perp.
\]
The six core--tail edges supply the displayed ambient frames.  Pulling back
the six tail--tail edges gives five identity maps and one $-J_c$, proving
both exhaustiveness and $J_cM_c=M_c$.  This formulation includes non-graph
planes; $M_c=\operatorname{ran}(I,S_c)$ with $S_c^2=-I$ is only a transverse
chart.
At $r=1$, $M_c=\operatorname{span}(1,\sigma i)$ and the six displayed frames
recover Lalonde's Theorem~4.5 normal form with sign $b=\sigma$.

\textbf{Cross-color sector packing.}
Choose a core unitary $W_c:K_c^3\to Q_c$.  For $c\ne d$, same-vertex
orthogonality at vertices $1,\ldots,9$ forces
\[
 W_d^*W_c=\begin{pmatrix}0&X&Y\\-X&0&Z\\-Y&-Z&0\end{pmatrix}.
\]
The lower blocks are negatives, not adjoints, because the matrix maps the
$c$ coordinate space to the $d$ coordinate space.  With
$\Omega_L=\left(\begin{smallmatrix}0&L\\-L&0\end{smallmatrix}\right)$, the
six tail compressions are
\[
 \Omega_Y,\ \Omega_Z,\ -\Omega_X,\
 \Omega_{Y+Z},\ \Omega_{Y-X},\ \Omega_{Z-X}.
\]
Their invertible coefficient transform implies, for $L=X,Y,Z$,
\[
 \Omega_LM_c\subseteq M_d^\perp,
 \qquad \Omega_LM_c^\perp\subseteq M_d.
\]

Since $M_c$ reduces $J_c$, write
\[
 M_c\cap\ker(J_c-\sigma iI)
   =\{(x,\sigma ix):x\in E_{c,\sigma}\},
 \quad e_{c,+}+e_{c,-}=r.
\]
Each $\Omega_L$ intertwines $J_c$ and $J_d$, so the plane flip says
$L E_{c,\sigma}\subseteq E_{d,\sigma}^{\perp}$.  For fixed $\sigma$, the
$n$ physical spaces
\[
 W_c(E_{c,\sigma}\oplus E_{c,\sigma}\oplus E_{c,\sigma})
\]
are pairwise orthogonal in dimension $nr$.  Hence
\[
 3\sum_c e_{c,+}\le nr,\qquad3\sum_c e_{c,-}\le nr.
\]
Adding gives $3nr\le2nr$, impossible for $r>0$.  Thus no quantum
$n$-coloring exists.  The explicit classical four-coloring of $G_{19}$ and
fresh colors on the joined clique give the matching $(n+1)$ upper bound.

\textbf{Verification.}
Two standard-library exact obstruction verifiers and a separate graph checker
accompany the proof.  They decode the graph6 checksum, verify the published
four-coloring and exhaustive non-three-colorability, replay independent
rational noncommutative SOS expansions, check the Walsh and skew-block
systems, multiply all six tail compressions, invert their coefficient
transform, and check the symbolic all-$n$ dimension gap.  No floating-point
computation enters the theorem.

\textbf{Scope and relation to prior work.}
The known rank-one theorem motivates the sign frame, but it is not used as a
reduction: the SOS proves the module-valued core directly, and the invariant
plane analysis handles genuinely higher-rank tails.  The contradiction
excludes every finite-dimensional representation of the $n$-coloring
relations, reducible or irreducible.  No finite local dimension supports $n$
colors.  At $n+1$ colors, the classical solution lifts to fixed dimension $d$
by tensoring with $I_d$ and to fixed rank $s$ by taking the direct sum of its
cyclic color relabelings and tensoring with $I_s$.  No claim is made about
infinite-dimensional or commuting-operator colorings.

\end{document}
